\documentclass[11pt]{article}
\usepackage[margin=1in]{geometry}
\usepackage{amsmath,amssymb,amsthm,mathtools}
\usepackage{hyperref}
\usepackage{booktabs}

\newtheorem{theorem}{Theorem}
\newtheorem{lemma}[theorem]{Lemma}
\newtheorem{proposition}[theorem]{Proposition}
\theoremstyle{definition}
\newtheorem{definition}[theorem]{Definition}
\theoremstyle{remark}
\newtheorem*{remark}{Remark}

\newcommand{\Z}{\mathbb{Z}}
\newcommand{\lowbits}{\operatorname{LowBits}}
\newcommand{\lean}[1]{\texttt{#1}}
\renewcommand{\Pr}{\mathop{\mathrm{Pr}}}

\title{Reusing the nonce in Dilithium's signing loop\\
\large A walk through the Lean proof \lean{Bathwater.lean}}
\author{}\date{}

\begin{document}
\maketitle

\begin{center}
\fbox{\parbox{0.9\textwidth}{\small\textbf{Disclaimer.} This document, the Lean proof and the
numerical experiments were produced mostly by Claude Opus~5.5 (Anthropic), under human
direction. The Lean statements are machine-checked; everything else, including the prose and the
experiments, has not been independently verified.}}
\end{center}

\begin{abstract}
Sprenkels and Westerbaan (``Don't throw your nonces out with the bathwater'', ePrint 2020/1158)
and Sprenkels' thesis (chapter~8) propose to keep part of the nonce $y=(y_1,\dots,y_\ell)$ of
Dilithium/ML-DSA across rejected signing attempts. We check these proposals in Lean~4 + Mathlib
(no \lean{sorry}, standard axioms only), in the model where every attempt gets a fresh uniform
challenge (the \emph{Trans} hybrid of the security proof). The thesis's \emph{alt-z} (z-check
first, polynomial by polynomial; on a failure at $y_i$ resample $y_1,\dots,y_i$) is
\emph{exact}: in every attempt, the joint law of the attempt number and the output equals that of
vanilla signing, so the vanilla zero-knowledge simulator applies unchanged. The ePrint's main
proposal \emph{alt-r0} (r$_0$-check first; on failure resample only $y_1$) is \emph{not}
perfectly zero-knowledge: on a Dilithium-shaped toy instance, two secret keys with the same public
key output the same signature with probabilities $1/6$ and $1/3$. Two other orderings fail the
same way. For the real parameters alt-r0 is nevertheless \emph{statistically} close to vanilla:
for all but a $2^{-127}$ fraction of matrices $A$, the alt-r0 signature distribution is within
$2^{-125}$ of vanilla's in $\ell_1$ for ML-DSA-44 ($2^{-138}$ and $2^{-155}$ for ML-DSA-65/87, with
bad fractions $2^{-139}$ and $2^{-157}$), per signature. The proof bounds the second moment, over $A$,
of the r$_0$ pass probability with the tail of $y$ fixed, by a Fourier expansion of the pass set
and the uncertainty-principle counting of \lean{RejOrder.lean}.
\end{abstract}

\tableofcontents

\section{The headline}\label{sec:headline}

\begin{center}\small
\begin{tabular}{@{}lll@{}}
\toprule
Policy & Source & Verdict\\
\midrule
vanilla, any check order & FIPS 204 & exact (by definition)\\
alt-z: z first, prefix resample & thesis Alg.~8.2 & \textbf{exact} (\lean{altZ\_exact})\\
alt-r0: r$_0$ first, keep $y_2..y_\ell$ & ePrint Alg.~2, thesis 8.3 &
  \textbf{leaks} on toy (\lean{toy\_altR0\_leaks});\\
 & & \textbf{close} for ML-DSA (\lean{altR0\_mldsa44})\\
r$_0$ first, prefix resample on z & ePrint Alg.~3 order &
  leaks on toy (\lean{toy\_r0FirstPrefixZ\_leaks})\\
Alg.~3: also head-only on r$_0$ & ePrint Alg.~3 & \textbf{close} for ML-DSA (\lean{alg3\_mldsa44})\\
z first, then keep $y_2..y_\ell$ on r$_0$ & (variant) &
  leaks on toy (\lean{toy\_zFirstThenY1\_leaks})\\
\bottomrule
\end{tabular}
\end{center}

\begin{itemize}
\item \textbf{Exact} means: for every attempt number $k$ and every $(y,c)$, the probability that
  attempt $k$ is the accepting one and outputs $(y,c)$ is the same as for vanilla signing. The
  output law, the law of the number of attempts, and their joint law are all vanilla's.
\item \textbf{Leaks} means: there are two secret keys with the \emph{same} public key, and a
  signature $(c,z)$, whose probability of being output differs between the keys. No simulator
  that sees only the public key can then be perfect.
\item The toy has $n=1$. It shows that the property the ePrint claims (perfect zero-knowledge) is
  false in general. It does \emph{not} show a practical leak for ML-DSA: there alt-r0 is provably
  close to vanilla (Section~\ref{sec:altr0stat}).
\item \textbf{Close} means: for all but a small fraction of matrices $A$, the alt-r0 signature law
  is within a small $\ell_1$ distance of vanilla's (table in Section~\ref{sec:altr0stat}).
\item The ePrint's \S4 argues that alt-r0 is perfectly zero-knowledge. Its simulator keeps
  $z_2,\dots,z_\ell$ and resamples $z_1$ and $c$. That is not what the scheme does: the scheme
  keeps $y_2,\dots,y_\ell$, so the tail of $z = y + c s_1$ changes with every new $c$. The thesis
  (\S8.A) withdraws the claim and gives a heuristic bound instead.
\end{itemize}

\section{Setting}\label{sec:setting}

A signing attempt draws $y=(y_1,\dots,y_\ell)$ uniformly from $(-\gamma_1,\gamma_1]^{n\ell}$,
computes the challenge $c$ from $w_1 = \mathrm{HighBits}(Ay)$, sets $z = y + c s_1$ and runs two
checks:
\begin{align*}
\text{z-check:}\quad & \|z_i\|_\infty < \gamma_1-\beta \text{ for every } i
  \qquad\text{(polynomial $i$ depends only on $y_i$ and $c$)},\\
\text{r$_0$-check:}\quad & \|\lowbits(Ay - c s_2)\|_\infty < \gamma_2 - \beta
  \qquad\text{(depends on all of $y$)}.
\end{align*}
Vanilla resamples all of $y$ after any failure. The proposals keep some polynomials of $y$.

\paragraph{The model (Trans hybrid).} The security proof of Fiat--Shamir with aborts
(\cite{BBDD23}, adapted in thesis \S8.4) first replaces the random-oracle call $c = H(M\|w_1)$ by
a fresh uniform $c$ in every attempt, and programs $H$ only for the accepted attempt. The
distance from the real scheme is bounded by the min-entropy of $w_1$ (the thesis does this for
alt-z). After that hop, zero-knowledge is a purely combinatorial question about the loop: what is
the law of the accepted $(y,c)$? That question is what the Lean file answers. Throughout:

\begin{description}
\item[A1] Each attempt's challenge $c$ is uniform on a finite set $C$, independent of everything
  else. The adversary sees only the accepted $(y,c)$; equivalently, the signature
  $(c, z=y+cs_1)$, from which $y$ is determined.
\item[A2] The z-check of polynomial $i$ depends only on $(y_i,c)$; the r$_0$-check may depend on
  all of $(y,c)$. Both are arbitrary predicates in the general theorems.
\item[A3] Resampled polynomials are fresh and uniform.
\end{description}
Not modelled: the Sign$\to$Trans hop itself (random-oracle programming and the min-entropy of
$w_1$ with a reused tail), the quantum random oracle, and timing.

\section{The loop as a chain}\label{sec:loop}

\begin{definition}[\lean{Loop}]
A rejection loop is given by an initial law $\pi$ on states, a sub-stochastic matrix
$B(s,s')$ (reject in state $s$, continue in $s'$) and $A(s,x)$ (accept in state $s$ with output
$x$). The state at attempt $k$ has sub-law $\mu_k = \pi B^k$ (\lean{state}), attempt $k$ accepts
with output $x$ with probability $\mathrm{out}_k(x) = \sum_s \mu_k(s) A(s,x)$ (\lean{outAt}), and
$\mathrm{out}_{<K}(x) = \sum_{k<K} \mathrm{out}_k(x)$ (\lean{outUpTo}).
\end{definition}

\begin{theorem}[\lean{Loop.tendsto\_outUpTo}]
Suppose $B, A \ge 0$, $A \le 1$, every row of $B$ sums to at most $1$, and for some $m \ge 1$ and
$\varepsilon > 0$, $m$ consecutive rejections have probability at most $1-\varepsilon$ from every
state. If $\nu$ solves $\nu = \pi + \nu B$, then
$\mathrm{out}_{<K}(x) \to \sum_s \nu(s) A(s,x)$ as $K\to\infty$.
\end{theorem}

$\nu(s)$ is the expected number of visits to $s$. The proof writes
$\nu = \sum_{k<K}\mu_k + \nu B^K$ and shows that the mass of $\nu B^K$ shrinks by $1-\varepsilon$
every $m$ steps. We need $m>1$ because some states cannot accept in one attempt: a kept
polynomial with a coefficient $\gamma_1$ fails the z-check for every $c$.

\begin{definition}[\lean{ReusePolicy}, \lean{ReusePolicy.loop}]
The nonce is $y : \iota \to V$ (here $\iota = \{1,\dots,\ell\}$, $V$ the values of one
polynomial). A reuse policy maps $(y,c)$ to either \emph{accept}, or \emph{reject keeping}
a set $K\subseteq\iota$. The state of its loop is the full nonce at the start of an attempt; the
first one is uniform. After a reject keeping $K$ the next nonce agrees with $y$ on $K$ and is
uniform elsewhere:
\[
  B(y,y') = \frac1{|C|}\sum_{c} [\text{$(y,c)$ rejects keeping $K$}]\,
    \frac{[\,y'|_K = y|_K\,]}{|V|^{|\iota\setminus K|}},\qquad
  A(y,(y,c)) = \frac{[\text{$(y,c)$ accepts}]}{|C|}.
\]
\end{definition}
Drawing the fresh polynomials at the end of the previous attempt rather than at the start of the
next one does not change anything; it keeps the state space simple.

\section{When reuse is exact}\label{sec:exact}

\begin{definition}[\lean{ReusePolicy.Invariant}]
A policy is \emph{invariant} if every reject event depends only on the polynomials it resamples:
if $(y,c)$ rejects keeping $K$, and $y'$ agrees with $y$ outside $K$, then $(y',c)$ also rejects
keeping $K$.
\end{definition}

\begin{theorem}[\lean{exact\_of\_invariant}, \lean{outAt\_eq\_of\_invariant}]
Let $p$ be the probability that one attempt with uniform $y$ and $c$ accepts. For an invariant
policy, $\mu_k$ is $(1-p)^k$ times the uniform law, for every $k$. Hence two invariant policies
that accept on the same $(y,c)$ have the same $\mathrm{out}_k$ for every $k$.
\end{theorem}

\begin{proof}[Idea]
It suffices that every column of $B$ sums to $1-p$: then uniform stays uniform, scaled by the
rejection probability. Fix $c$ and $K$, and let $f(y)$ be the indicator that $(y,c)$ rejects
keeping $K$. By invariance $f$ ignores the coordinates in $K$. So among the $y$ that agree with a
given $y'$ on $K$, the fraction with $f(y)=1$ equals the overall fraction (\lean{sum\_agree},
via $y \leftrightarrow (y|_K, y|_{\iota\setminus K})$). Hence
$\sum_y f(y)\,[y|_K = y'|_K]/|V|^{|\iota\setminus K|} = \Pr_y[f(y)=1]$. Summing over $K$ and $c$
gives $1-p$.
\end{proof}

Vanilla (always keep $\emptyset$) is trivially invariant (\lean{vanilla\_invariant}).

\begin{theorem}[\lean{altZ\_exact}]
For any z-checks satisfying A2 and any r$_0$-check, alt-z has the same $\mathrm{out}_k$ as
vanilla for every $k$.
\end{theorem}
Alt-z keeps $\{j : j > i\}$ when $i$ is the first polynomial that fails its z-check, and keeps
nothing when the z-check passes and the r$_0$-check fails. The event ``$i$ is the first failure''
depends only on $y_1,\dots,y_i$ and $c$ (\lean{firstFail\_congr}), and these are exactly the
resampled polynomials. So alt-z is invariant (\lean{altZ\_invariant}); it accepts exactly when
vanilla does (\lean{altZ\_accept}).

\begin{remark}
This explains which orderings work. Before a reject keeps $y_j$, \emph{no} check that has been
looked at may depend on $y_j$. The r$_0$-check depends on all of $y$, so any policy that keeps
something after the r$_0$-check has been evaluated breaks invariance. This covers: alt-r0; the
r$_0$-first order of ePrint Algorithm~3 (whose kept tail has passed an r$_0$-check); and
``z first, keep the tail on an r$_0$ failure''. Keeping $y_1$ after a failure at $y_2$ breaks it
too (the ePrint and thesis both note this).
\end{remark}

\section{Approximate invariance}\label{sec:approx}

Exactness (Section~\ref{sec:exact}) rests on one fact: if the nonce is uniform at the start of an
attempt and the attempt rejects, the nonce is again uniform, with mass $1-p$. For policies where
this holds only approximately, the error does not compound. Let $\pi$ be uniform and
\[
  \mathrm{defect}(P) = \big\|\pi B - (1-p)\pi\big\|_1
  \qquad(\lean{ReusePolicy.defect}),
\]
the $\ell_1$ distance between the state after one rejected attempt from a uniform nonce and the
exact $(1-p)\pi$. Invariant policies have defect $0$ (\lean{defect\_of\_invariant}).

\begin{theorem}[\lean{state\_close}, \lean{outUpTo\_close} etc.]
Let $p = \Pr[\text{a fresh attempt accepts}] > 0$. For every $k$,
\[
  \big\|\mathrm{state}_k - (1-p)^k\pi\big\|_1 \le \mathrm{defect}\cdot\sum_{j<k}(1-p)^j
  \le \mathrm{defect}/p.
\]
If $Q$ is invariant and accepts on the same $(y,c)$ as $P$, then in every attempt the joint law
of ``attempt $k$ accepts with output $x$'' is within $\mathrm{defect}/p$ of $Q$'s, within the first
$K$ attempts within $K\cdot\mathrm{defect}/p$, and $P$ is still running after $K$ attempts with
probability at most $(1-p)^K + \mathrm{defect}/p$.
\end{theorem}

The proof is one line: with $e_k = \mathrm{state}_k - (1-p)^k\pi$,
$e_{k+1} = e_k B + (1-p)^k(\pi B - (1-p)\pi)$, and $B$ does not increase the $\ell_1$ norm
(rows of $B$ sum to at most $1$). No limits or fixed points are needed, and the statements cover
the joint law of attempt number and output. This is the tool for ePrint Algorithm~3
(r$_0$ first; resample only $y_1$ on an r$_0$ failure, a prefix on a z failure), whose output law
has no closed form.

\section{Alt-r0: the output law}\label{sec:altr0}

In alt-r0 the polynomials of $y$ outside a set $H$ (in the papers $H = \{1\}$) are kept while the
r$_0$-check fails; if it passes and the z-check fails, everything is resampled. Write $\rho(y)$
for the probability, over fresh $y_H$ and $c$, that the r$_0$-check passes; it depends only on
the kept part $t = y_{\bar H}$ (\lean{AltR0Set.rho}). Call $y$ \emph{stuck} if $\rho(y) = 0$: then
the loop never ends.

\begin{theorem}[\lean{altR0Set\_law}, \lean{vanilla\_law}]
Let $P$ be the acceptance probability of one fresh attempt (\lean{pacc}), and assume $P > 0$. The
probability that alt-r0 outputs $(y,c)$ (within $K$ attempts, $K\to\infty$) tends to
\[
  U(y,c) = \frac{\kappa\,[(y,c)\text{ passes}]\,[\rho(y)\neq 0]}{|Y|\,|C|\,\rho(y)},
  \qquad \kappa = \frac1{1-F},\quad
  F = \mathop{\mathbb E}_y\Big[\frac{\Pr_c[\text{r}_0 \text{ passes, z fails}]}{\rho(y)};\ \rho(y)\ne0\Big],
\]
whereas vanilla's law is $V(y,c) = [(y,c)\text{ passes}]/(|Y||C|P)$.
\end{theorem}

Intuitively, a tail that rarely passes the r$_0$-check is stuck for many attempts and gets many
draws of $(y_H, c)$. Since $t = z_{\bar H} - c\,s_{1,\bar H}$ and $\rho$ depends on $A$ and $s_2$,
the weight $1/\rho$ depends on the secret; Section~\ref{sec:toy} turns this into a leak on a toy.
The proof needs loop semantics that allow stuck states: \lean{Loop.tendsto\_of\_fixed\_off}
computes the output law from a nonnegative solution of $\nu = \pi + \nu B$ off a closed set of
non-accepting states (here $\nu(y) = \kappa/(|Y|\rho(y))$), by treating entry into that set as an
extra output.

\begin{theorem}[\lean{altR0Set\_close}]
For $\mu_0 > 0$ let $D = \mathop{\mathbb E}_y |\rho(y) - \mu_0| / \mu_0$. If $\mu_0 D < P$ then
\[
  \sum_x |U(x) - V(x)| \le \frac{2\mu_0 D}{P - \mu_0 D},
  \qquad 1 - \sum_x U(x) \le \frac{\mu_0 D}{P - \mu_0 D}.
\]
\end{theorem}
The second quantity is the probability that alt-r0 never terminates. The proof groups $y$ by its
kept part: within a class $\rho$ is constant, the class average of $\Pr_c[\text{pass}]$ is at most
$\rho$, and stuck classes never pass.

So alt-r0 is exactly as close to vanilla as $\rho$ is to a constant, on average over tails.

\section{Alt-r0 for ML-DSA: statistically close}\label{sec:altr0stat}

\subsection{The result}

Take $\mu$ to be the probability that a \emph{uniform} $w$ passes the r$_0$-check,
$\mu = (\#\text{passing values}/q)^{kn}$, and let $\rho_c(t)$ be the probability, over $y_{h_0}$
uniform in its box, that $A\cdot(y_{h_0}, t)$ passes the r$_0$-check with shift $c s_2$
(\lean{AltR0.rhoC}).

\begin{theorem}[\lean{AltR0.dev\_mldsa44\_sharp}, \lean{\_65\_sharp}, \lean{\_87\_sharp}]
For every $c$ with $\|c s_2\|_\infty \le \beta$, every choice of the resampled polynomial $h_0$,
and $y_{h_0}$ ranging over either the full box or the z-passing box shifted by $c s_1$,
\[
  \mathop{\mathbb E}_A \mathop{\mathbb E}_t |\rho_c(t) - \mu| \le \mu\cdot 2^{-147},\
  \mu \cdot 2^{-182},\ \mu\cdot 2^{-190}
\]
for ML-DSA-44, -65, -87, with $A$ uniform in $R_q^{k\times\ell}$ and $t$ uniform.
\end{theorem}

The second moment behind it is $\mathop{\mathbb E}_A \mathop{\mathbb E}_t (\rho_c(t) - \mu)^2
\le \mu^2 2^{-294}, \mu^2 2^{-364}, \mu^2 2^{-380}$ (\lean{var\_mldsa44\_sharp} etc.), and
with the $1/q$ slot bound of Section~\ref{sec:slotq} it is
$\le \mu^2 2^{-382}, \mu^2 2^{-420}, \mu^2 2^{-472}$ (\lean{var\_mldsa44\_q} etc.). Combining
with Section~\ref{sec:altr0} (the full box controls $D$, the z-passing box gives
$P \ge \mu p_z - \text{error}$ with $p_z \approx 0.54, 0.62, 0.66$ the z pass probability) and
Markov's inequality over $A$ applied to that \emph{second} moment (a matrix is bad if the mean
squared deviation exceeds $(\mu\varepsilon)^2$; for the others the mean deviation is below
$\mu\varepsilon$ by Cauchy--Schwarz, \lean{main2}):

\begin{theorem}[\lean{AltR0Final.altR0\_mldsa44}, \lean{\_65}, \lean{\_87}]
For any finite challenge set $C$ and $c s_1$, $c s_2$ bounded by $\beta$, and any $h_0$:
\begin{center}
\begin{tabular}{@{}lccc@{}}
\toprule
 & bad fraction of $A$ & $\sum_x|U(x)-V(x)|$ & $\Pr[\text{never terminates}]$\\
\midrule
ML-DSA-44 & $\le 2^{-127}$ & $\le 2^{-125}$ & $\le 2^{-126}$\\
ML-DSA-65 & $\le 2^{-139}$ & $\le 2^{-138}$ & $\le 2^{-139}$\\
ML-DSA-87 & $\le 2^{-157}$ & $\le 2^{-155}$ & $\le 2^{-156}$\\
\bottomrule
\end{tabular}
\end{center}
The last two columns hold for every $A$ outside the bad set. The bad fraction scales as
$2\cdot 2^{-k}/\varepsilon^2$ ($k = 382, 420, 472$) and the distance as $4\varepsilon$, so the
trade-off can be moved; e.g.\ keeping the distance at $2^{-95}$ for ML-DSA-44 makes the bad
fraction $2^{-187}$. (Before the $1/q$ slot bound, $k = 294, 364, 380$ gave
$2^{-97}/2^{-96}$, $2^{-121}/2^{-119}$, $2^{-127}/2^{-124}$.) (Applying
Markov to the first moment instead, \lean{altR0\_mldsa44\_eps}, gives only $2^{-73}$ and
$2^{-71}$.)
\end{theorem}

These are per-signature distances in the Trans model; over $N$ signatures they add up ($N$ times).
They are crude: the thesis heuristic suggests about $2^{-2289}$.

\paragraph{Per attempt (section \lean{AltR0JointPart}).} The theorem above is about the output law
in the limit. Section~\ref{sec:approx} gives more: alt-r0's defect is at most $2X(A)$
(\lean{altR0Set\_defect\_le}; the column sum at $y'$ is $1-\rho(\text{tail of }y')$ plus a constant
from z failures, which resample everything), so for the same good matrices and every $k$
(\lean{altR0\_joint\_mldsa44}, \lean{\_65}, \lean{\_87}):
the joint law of (accepting attempt $<k$, output) is within $k\cdot 2^{-125}$, $k\cdot 2^{-138}$,
$k\cdot 2^{-155}$ of vanilla's, and alt-r0 is still running after $k$ attempts with probability
at most $(1-\mu/2)^k + 2^{-125}$ (resp.\ $2^{-138}$, $2^{-155}$). This covers the number of
attempts, which is what timing reveals, and justifies an attempt cap such as CIRCL's $576$.

\subsection{Why it works: the second moment over $A$}

Fix $c$ (so the shift $v = c s_2$) and write $G \subseteq \Z_q^{kn}$ for the set of $w$ that pass the
r$_0$-check; $G$ is a product of per-coordinate sets of density $1 - p_R$. Then
\[
  \mathop{\mathbb E}_A \mathop{\mathbb E}_t (\rho_c(t) - \mu)^2
  = \mathop{\mathbb E}_{y_{h_0}, y'_{h_0}, t}\ \mathop{\mathbb E}_A
    \big[(1_G(Ay) - \mu)(1_G(Ay') - \mu)\big],
  \qquad y = (y_{h_0}, t),\ y' = (y'_{h_0}, t),
\]
a sum over pairs of nonces that \emph{share the tail}.

\paragraph{Fourier (section \lean{FourierPart}, \lean{cov\_le}).} For $A$ uniform, $(Ay, Ay')$ is uniform on
a subgroup, and
\[
  \Big|\mathop{\mathbb E}_A\big[(1_G(Ay) - \mu)(1_G(Ay') - \mu)\big]\Big|
  \le \mu^2 \!\!\sum_{\substack{\xi,\xi' \neq 0\\ \xi\cdot Ay + \xi'\cdot Ay' \equiv 0}}\!\!
      b^{\mathrm{wt}(\xi) + \mathrm{wt}(\xi')},
  \qquad b = \frac{p_R}{1-p_R} \approx 2^{-10.25}.
\]
The point is that $1_G$ is a product: its Fourier coefficient at $\xi$ is a product over
coordinates, and every nonzero coordinate costs a factor $b$ (the bad set is small). This
replaces the Bonferroni truncation of \lean{RejOrder.lean}, which with only one fresh polynomial
gives just $2^{-15}$.

\paragraph{Support (\lean{degS\_sub}, \lean{dualIn\_of\_orth}).} An orthogonal pair $(\xi,\xi')$ has
its dual (NTT) coordinates on the \emph{degenerate slots} of $(y,y')$. For a pair sharing the tail,
a slot is degenerate only if $y_{h_0} - y'_{h_0}$ vanishes there or the whole tail vanishes there.
Each happens with probability at most $1/w$ resp.\ $(2\gamma_1)^{-(\ell-1)}$ by the window argument
($w$ the width of the head box), so $d$ given slots are all degenerate with probability at most
$p^d$, $p = 1/w + (2\gamma_1)^{-(\ell-1)}$ (\lean{prob\_Z}, \lean{prob\_deg\_ge}).

\paragraph{Counting (\lean{wsum\_le}, \lean{wcount\_le\_EN}).} The $b$-weighted number of test
vectors on $d$ slots is bounded by the gap counts of \lean{RejOrder.lean}: a test vector supported
on coefficient set $T$ with duals on $d$ slots is one of at most $q^{\mathrm{xg}(d,T)}$, and
$\sum_{|T| = t} q^{\mathrm{xg}}$ is evaluated with \lean{coefK}. Pairs with more than
$d_m$ degenerate slots are charged the trivial bound $1$.

\paragraph{Pairing $\xi'$ with $\xi$ (section \lean{AltR0SharpPart}).} The covariance bound above weights
$\xi$ and $\xi'$ independently. In fact, at every degenerate slot the pair is tied:
where the tail does not vanish, orthogonality forces $\tilde\xi'(k) = -\tilde\xi(k)$; where the
tail vanishes and both heads are nonzero, $\tilde\xi'(k) = -(\hat y_{h_0}(k)/\hat y'_{h_0}(k))
\tilde\xi(k)$. So unless a degenerate slot is \emph{exceptional} (tail and one head vanish there),
$\xi' = \Phi\xi$ for a bijection $\Phi$ that rescales dual coordinates slot by slot
(\lean{orth\_eq\_pairMap}). With $b^{\mathrm{wt}\,\xi + \mathrm{wt}\,\Phi\xi} \le
(b^{2\mathrm{wt}\,\xi} + b^{2\mathrm{wt}\,\Phi\xi})/2$ and $\Phi$ permuting the test vectors on the
degenerate slots, the pair sum is at most the $b^2$-weighted count of single test vectors
(\lean{pair\_sum\_le}): squared weight, and no square of the count. Exceptional slots occur with
probability about $2(2\gamma_1)^{-\ell}$ each (section \lean{AltR0ExcPart}, \lean{prob\_deg\_exc});
triples with one keep the crude bound.

\paragraph{Numbers.} The crude version gives (\lean{var\_le}, \lean{var\_rel})
\[
  \mathop{\mathbb E}_A\mathop{\mathbb E}_t (\rho_c(t)-\mu)^2
  \le \mu^2 \sum_{d=1}^{d_m}\binom nd p^d E_b(d)^2 + \binom{n}{d_m+1}p^{d_m+1},
\]
with $E_b(d)$ the bound on the weighted count. Cleared of denominators, the inequality
``$\le \mu^2/2^{k}$'' is an integer inequality checked by \lean{decide +kernel}, with $k = 246,
290, 266$ ($d_m = 19, 19, 17$; $t$ up to $250, 400, 500$ in the gap counts; about $10$\,s each).
Cauchy--Schwarz halves the exponent (\lean{dev\_rel}): about $2^{-123}, 2^{-145}, 2^{-133}$
(these crude instances are no longer kept in the file). The
sharp version (\lean{var\_rel\_sharp}, \lean{dev\_rel\_sharp}) has two sums, the main one with
$b^2$-weighted counts up to $d_1 = 25, 26, 27$ degenerate slots and the exceptional one as before,
and is checked with $k = 294, 364, 380$; hence $2^{-147}, 2^{-182}, 2^{-190}$. Here the weighted
counts are bounded by rows (\lean{wcount\_le\_rows}) rather than by truncating in $|T|$.

\subsection{A slot condition costs $1/q$}\label{sec:slotq}

The window argument bounds the probability that a box polynomial has prescribed NTT values at
$r$ given slots by $w^{-r}$, $w$ the box width ($2^{18}$ or $2^{20}$). The truth is about
$q^{-r}$ ($q \approx 2^{23}$). Section \lean{PackingPart} proves
\[
  \#\{f \in \text{box of width } w : \hat f|_S = \tau\} \cdot q^{|S|} \le (w + 2nB)^n
  \qquad\text{whenever } (B+1)^n > q^{|S|}
\]
(\lean{Packing.card\_coset\_le}), by packing:
\begin{enumerate}
\item Pigeonhole on $[0,B]^n$ gives a nonzero $v$ with $|v_s| \le B$ and $\hat v|_S = 0$.
\item The multiples of $v$ in $R = \Z[X]/(X^n+1)$ keep $\hat{\ }|_S = 0$. Since $X^{256}+1$ is the
  $512$-th cyclotomic polynomial, it is irreducible over $\mathbb Q$, so the negacyclic matrix of
  $v$ is invertible over $\mathbb Q$ (section \lean{NcInvPart}, \lean{ncMat\_det\_ne\_zero}).
  Every integer $x$ is then $x = v\cdot\lfloor v^{-1}x\rfloor + \mathrm{red}(x)$ with
  $|\mathrm{red}(x)_s| \le nB$ and the same values on $S$.
\item Reducing the $q^{|S|}$ vectors of $[0,q)^{|S|}\times 0$ (distinct on $S$ by the window
  argument) gives $q^{|S|}$ short representatives; translating the solutions by them packs
  $q^{|S|}$ disjoint copies into a box of width $w + 2nB$.
\end{enumerate}
With $B = 5$ (valid up to $28$ slots, $6^{256} > q^{28}$) the constant $((w+2nB)/w)^n$ is about
$2^{3.6}$ (ML-DSA-44) and $2^{0.9}$ (65/87). Section \lean{AltR0QPart} redoes the slot
probabilities with $p = 1/q + (2\gamma_1)^{-(\ell-1)}$: \lean{prob\_deg\_ge\_q} (constant once) and
\lean{prob\_deg\_exc\_q} (constant squared: two head conditions). \lean{var\_rel\_gen} takes
these probabilities as hypotheses; the kernel checks (\lean{numQ\_44\_full} etc.) use
$d_1 = 27$, $d_2 = 19$ and give $k = 382, 420, 472$.

\subsection{Where the bound is loose}

\begin{itemize}
\item Triples with an exceptional slot still use the crude bound, and they dominate: with the
  $1/q$ bound the main part would give about $2^{-517}$ for ML-DSA-44, the exceptional part
  $2^{-386}$.
\item Going from the second moment to $\mathop{\mathbb E}|\rho - \mu|$ halves the exponent.
\end{itemize}

\section{ePrint Algorithm 3 for ML-DSA}\label{sec:alg3}

Algorithm~3 checks r$_0$ first; on an r$_0$ failure it resamples only the head $y_1$, and on a
z failure at polynomial $i$ (the first failing one) it resamples $y_1,\dots,y_i$ and keeps the
rest (\lean{alg3}; it accepts exactly when vanilla does, \lean{alg3\_accept}). Its output law has
no closed form, so we use Section~\ref{sec:approx}.

\paragraph{The defect.} With the head the first polynomial, an r$_0$ failure and ``z first fails
at the head'' keep the same set, so together they move to the head fibre with probability
$1 - q_h\,\rho_z(\text{tail})$, $q_h = (zw/2\gamma_1)^n$ the probability that the head passes z
and $\rho_z$ the r$_0$ pass probability with the head in the z-passing box. The full-box $\rho$
cancels. Each later failure point $i$ contributes a term with $y_2,\dots,y_i$ fresh, which after
summing over the kept part is again a deviation of $\rho_z$ over a uniform tail. Centring (the
column sums have mean exactly $1-p$) costs a factor $2$:
\[
  \mathrm{defect}(\text{Alg.\,3}) \le 2\ell\cdot X_z(A)
  \qquad(\lean{alg3\_defect\_le}),
\]
with $X_z = \mathop{\mathbb E}_c\mathop{\mathbb E}_t|\rho_{c}(t) - \mu|$ (z-passing head box) as in
Section~\ref{sec:altr0stat}. The acceptance probability is vanilla's, $\ge \mu(p_z - \varepsilon)$
for good $A$ (\lean{alg3\_pAcc\_eq}, \lean{pacc\_ge}).

\begin{theorem}[\lean{AltR0Final.alg3\_mldsa44}, \lean{\_65}, \lean{\_87}]
Outside the same bad set of matrices as for alt-r0 (fractions $2^{-127}$, $2^{-139}$, $2^{-157}$),
for every number of attempts $k$: the joint law of (accepting attempt $< k$, output) of
Algorithm~3 is within $k\cdot 2^{-123}$, $k\cdot 2^{-135}$, $k\cdot 2^{-152}$ of vanilla's in
$\ell_1$, and Algorithm~3 is still running after $k$ attempts with probability at most
$(1-\mu/2)^k + 2^{-123}$ (resp.\ $2^{-135}$, $2^{-152}$).
\end{theorem}

The factor $k$ is the price of the joint law; for the output alone, the cutoff $k$ can be taken
where $(1-\mu/2)^k$ is negligible (a few hundred attempts). The bad set reuses \lean{Bad2}, which
also asks for the full-box second moment; Algorithm~3 needs only the z-box one.

\section{The toy counterexamples}\label{sec:toy}

The instance (section \lean{Toy}): $q=17$, $n=1$ so $R_q = \Z/17$, $\gamma_2 = 2$
($2\gamma_2 \mid q-1$), $\gamma_1 = 3$, $\eta=\tau=1$ (so $c=\pm1$, $\beta = 1$), $\ell=2$,
$k=1$, $A = (4\ \ 7)$. $\lowbits$ is FIPS~204's \textsf{Decompose}, including its corner case.
The keys
\[
  s_1 = (-1,-1),\ s_2 = -1 \qquad\text{and}\qquad s_1 = (1,0),\ s_2 = 1
\]
have the same public key $t = As_1 + s_2 = 5$ (\lean{same\_pk}). The signature
$(c,z) = (1,(0,0))$ comes from $y=(1,1)$ under the first key and $y=(-1,0)$ under the second
(\lean{y1\_sig}, \lean{y2\_sig}). Its probability:

\begin{center}
\begin{tabular}{@{}lccl@{}}
\toprule
Policy & key 1 & key 2 & Lean\\
\midrule
vanilla & $1/4$ & $1/4$ & \lean{toy\_vanilla\_same}\\
alt-r0 & $1/6$ & $1/3$ & \lean{toy\_altR0\_leaks}\\
r$_0$ first, prefix z & $3/11$ & $5/22$ & \lean{toy\_r0FirstPrefixZ\_leaks}\\
z first, keep tail on r$_0$ & $4/19$ & $1/4$ & \lean{toy\_zFirstThenY1\_leaks}\\
\bottomrule
\end{tabular}
\end{center}

Each entry is a limit of $\mathrm{out}_{<K}$ as $K\to\infty$, proved with
Theorem~\lean{Loop.tendsto\_outUpTo}. The visit vector $\nu$ (36 rationals) was found by exact
linear algebra outside Lean. The kernel checks, by \lean{decide +kernel}, that $\nu = \pi + \nu B$,
that rows of $B$ sum to at most $1$, that three consecutive rejections have probability at most
$1 - 1/200$, and the value $\sum_s \nu(s)A(s,x)$. The policies are the same Lean definitions
(\lean{altR0}, \lean{r0FirstPrefixZ}, \lean{zFirstThenY1}, \lean{vanilla}) as in the general
theorems, instantiated with the toy checks.

On the full ePrint Algorithm~3 (alt-r0 together with prefix resampling after the r$_0$-check)
the toy has tails that can never be accepted: the tail is kept on an r$_0$ failure and on a z
failure at $y_1$, and for some tails one of the two always fails. This does not happen for real
parameters, but it is why Algorithm~3 has no toy theorem; its two ingredients each leak.

\section{Min-entropy of the commitment}\label{sec:minentropy}

The hop from real signing to the Trans model (fresh challenges) loses, per signing query, roughly
$q_H\,\varepsilon$ times the expected number of attempts, where $\varepsilon$ bounds the
probability of guessing an attempt's commitment $w_1 = \mathrm{HighBits}(Ay)$ given everything
else \cite{BBDD23}; the thesis (\S8.4) adapts this to nonce reuse, with a guessing probability
$\varepsilon_j$ per attempt. With reuse the tail of $y$ can be correlated with earlier attempts,
but in alt-r0 and alt-z the polynomial $y_{h_0}$ (the first one) is fresh in every attempt. So it
suffices to bound $\Pr_{y_{h_0}}[\mathrm{HighBits}(A\cdot(y_{h_0},t)) = W]$ for every fixed tail
$t$ and target $W$ (section \lean{MinEntropyPart}).

\begin{theorem}[\lean{guess\_le}, \lean{minEntropy}]
If the entry $A_{i,h_0}$ has $z$ zero NTT slots, then for every tail $t$ and every $W$
\[
  \Pr_{y_{h_0}}\big[\mathrm{HighBits}((A y)_i) = W_i\big]
  \le \Big(\frac{2\gamma_2 + 1}{2\gamma_1}\Big)^{n} (2\gamma_1)^{z}.
\]
\end{theorem}
The value $(Ay)_i$ is determined by its high and low bits (\lean{eq\_of\_bits}), and the low bits
take at most $2\gamma_2+1$ values per coefficient. Two heads with the same $(Ay)_i$ differ by an
element killed by $A_{i,h_0}$, whose NTT vanishes off the $z$ zero slots, so by the window
argument at most $(2\gamma_1)^z$ heads share a value. The thesis's argument is the case $z = 0$
(an invertible entry).

\begin{theorem}[\lean{card\_bad\_le}, \lean{minEntropy\_mldsa44}, \lean{\_65}, \lean{\_87}]
The matrices in which every entry of column $h_0$ has at least $d$ zero slots are a fraction at
most $(\binom nd / q^d)^k$. With $d = 2$:
\begin{center}
\begin{tabular}{@{}lcc@{}}
\toprule
 & guessing probability $\varepsilon$ & outside a fraction of $A$\\
\midrule
ML-DSA-44 & $\le 2^{-99}$ & $\le 2^{-124}$\\
ML-DSA-65 & $\le 2^{-236}$ & $\le 2^{-186}$\\
ML-DSA-87 & $\le 2^{-236}$ & $\le 2^{-248}$\\
\bottomrule
\end{tabular}
\end{center}
\end{theorem}
With $d = 1$ (the thesis's invertibility condition) the numbers are $\varepsilon \le 2^{-118}$
outside a $2^{-60}$ fraction for ML-DSA-44; $d = 2$ trades $18$ bits of $\varepsilon$ for $64$
bits of the fraction. The game-hopping argument itself (from these $\varepsilon_j$ to the loss in
the reduction) is that of \cite{BBDD23} and the thesis, and is not formalised here.

\section{What is not covered}\label{sec:notcovered}

\begin{itemize}
\item \textbf{Sign $\to$ Trans.} The min-entropy it needs is proved
  (Section~\ref{sec:minentropy}); the game hops themselves are not formalised.
\item \textbf{Many signatures.} The alt-r0 distances are per signature; there is no hybrid
  theorem over many signatures, and the numbers are far from what is believed true.
\item \textbf{Timing.} Exactness is about the accepted signature and the number of attempts. With
  alt-z the work in an attempt depends on how many polynomials are resampled, which reveals the
  index of the first failing z polynomial. That is the same kind of information as in
  \lean{RejOrder.lean} (which check failed, and where); it is not analysed here.
\item The toy has $n = 1$; there is no NTT structure.
\end{itemize}

\section{Map of the Lean file}\label{sec:map}

Everything is in \lean{Bathwater.lean}, which imports Mathlib and \lean{RejOrder.lean} (for the
module/NTT machinery and the gap counts). Its sections, in order:
\begin{center}
\begin{tabular}{@{}ll@{}}
\toprule
Section & Contents\\
\midrule
\lean{BathwaterCore} & loop semantics (\lean{Loop}), \lean{ReusePolicy}, \lean{exact\_of\_invariant},
  \lean{altZ\_exact};\\
 & alt-r0 output law and distance (\lean{AltR0Set}); toy counterexamples (\lean{Toy})\\
\lean{FourierPart} & covariance of product-set indicators (\lean{cov\_le})\\
\lean{AltR0Part} & shared-tail slots, weighted counts, \lean{var\_le}, \lean{dev\_rel}\\
\lean{PackingPart} & a slot condition costs $1/q$ (\lean{card\_coset\_le})\\
\lean{NcInvPart} & $X^{256}+1$ irreducible, so \lean{ncMat\_det\_ne\_zero}\\
\lean{AltR0ExcPart} & probability of exceptional slots (\lean{prob\_deg\_exc})\\
\lean{AltR0QPart} & slot probabilities with $1/q$ (\lean{prob\_deg\_ge\_q}, \dots)\\
\lean{AltR0SharpPart} & pairing map, \lean{pair\_sum\_le}, \lean{wcount\_le\_rows},
  \lean{var\_le\_sharp}, \lean{dev\_mldsa*\_sharp}\\
\lean{AltR0GenPart} & \lean{var\_rel\_gen}, \lean{var\_mldsa*\_q}\\
\lean{AltR0FinalPart} & instantiation for ML-DSA, \lean{altR0\_mldsa44/65/87}\\
\lean{AltR0JointPart} & alt-r0 per attempt, \lean{altR0\_joint\_mldsa*}\\
\lean{Alg3Part} & Algorithm 3: \lean{alg3\_defect\_le}, \lean{alg3\_mldsa*}\\
\lean{MinEntropyPart} & guessing the commitment with a fresh head, \lean{minEntropy\_mldsa*}\\
\bottomrule
\end{tabular}
\end{center}
The kernel checks bound the weighted count by rows: the sum over coordinate sets factors as a
$k$-th power of a row sum (\lean{wcount\_le\_rows}), so each check multiplies numbers of at most
about $50{,}000$ bits. \lean{build.sh} builds \lean{RejOrder.lean} and then \lean{Bathwater.lean}
(about $165$\,s and $12$\,GB for the latter).

\begin{thebibliography}{9}
\bibitem{SW20} A.~Sprenkels and B.~Westerbaan. Don't throw your nonces out with the bathwater:
  speeding up Dilithium by reusing the tail of $y$. Cryptology ePrint Archive 2020/1158.
\bibitem{Spr24} A.~Sprenkels. PhD thesis, Radboud University, 2024, chapter~8.
\bibitem{BBDD23} M.~Barbosa et al. Fixing and mechanizing the security proof of
  Fiat--Shamir with aborts and Dilithium. CRYPTO 2023.
\end{thebibliography}

\end{document}
